\documentclass[aspectratio=1610,onlymath]{beamer}
% \documentclass[aspectratio=1610,onlymath,handout]{beamer}

\input{macros-lecture}

\defineTitle{5}{DPLL und SMT}{30. April 2026}

% Resetting the definition of \\ is essential to avind global layout changes in beamer when loading minted.
% It does not seem to break other functionality in obvious ways.
% If the global override causes any issues, the more local one using \addtobeamertemplate{frametitle} (in document) might be a workaround.
\let\beameroriginaldoublebackslash\\
\usepackage{minted}
\let\\=\beameroriginaldoublebackslash

%%% Forests, vertical layouts, and beamer:
%%% https://tex.stackexchange.com/questions/5073/making-a-simple-directory-tree
%%% https://tex.stackexchange.com/questions/269734/how-to-adjust-large-number-of-nodes-in-forest-environment
\usepackage{forest}
\usetikzlibrary{overlay-beamer-styles}
\tikzset{
    invisible/.style={opacity=0,text opacity=0},
    visible on/.style={alt=#1{}{invisible}},
    alt/.code args={<#1>#2#3}{%
      \alt<#1>{\pgfkeysalso{#2}}{\pgfkeysalso{#3}} % \pgfkeysalso doesn't change the path
    }
}
\forestset{
  visible on/.style={
    for tree={
      /tikz/visible on={#1},
      edge+={/tikz/visible on={#1}}}},
  vel/.style={edge label={node[inner sep=1pt, inner xsep=3pt,midway,above,anchor=south west,xshift=0pt,font=\scriptsize]{#1}}},
  eqnode/.style={edge={draw=none}, edge label={node[inner sep=1pt, inner xsep=3pt,right,anchor=east,xshift=0pt]{$\equiv$}}}
}
\begin{document}

\maketitle

% \addtobeamertemplate{frametitle}{%
%   \let\\=\beameroriginaldoublebackslash
% }{}

\begin{frame}\frametitle{Wiederholung: Logisches Schließen}

\emph{Allgemeine Fragestellungen des logischen Schließens}
\begin{itemize}
\item Logische Folgerung: Gilt $\mathcal{F}\models G$?
\item Allgemeingültigkeit: Gilt $\models F$?
\item Unerfüllbarkeit: Gilt $\models \neg F$, d.h. $F\models\bot$?
\end{itemize}\bigskip

\emph{Einige einfache Methoden des logischen Schließens}
\begin{itemize}
\item Wahrheitswertetabelle
\item Erfüllbarkeitstest mit DNF
\item Widerlegbarkeitstest mit KNF
\item Erfüllbarkeitstest mit Resolution
\end{itemize}
$\leadsto$ Alle schlimmstenfalls exponentiell und oft zu ineffizient

\end{frame}

\sectionSlide{Die DPLL-Methode}

\begin{frame}\frametitle{Ein effizienteres Verfahren}

\emph{Idee:} Implementiere den Erfüllbarkeitstest als \redalert{Suche nach einem Modell}
\begin{itemize}
\item \alert{Teste systematisch Belegungen} für Atome
\item \alert{Bilde Ableitungen} um die Konsequenzen einer Belegung zu ermitteln
\item Erkenne Widersprüche und verwende \alert{Backtracking} um andere Belegungen zu \ghost{testen}
\end{itemize}
\medskip

DPLL arbeitet mit Formeln in \alert{Klauselform} (muss vorab gebildet werden)\\[1ex]
{\footnotesize \mbox{}~~~~~~Wir testen Erfüllbarkeit $\leadsto$ "`kleine"' (polynomielle) Klauselform mit Hilfsatomen zulässig

}\bigskip


\anybox{strongyellow}{\emph{Hintergrund:} DPLL steht für die Namen Davis, Putnam, Logemann, Loveland
\begin{itemize}
\item Davis \& Putnam hatten einen resolutionsbasierten Vorgängeralgorithmus entwickelt [1960]
\item Davis, Logemann \& Loveland entwickelten daraus die hier vorgestellte Version [1961]
\item Weitere Verbesserungen wurden erst viel später entwickelt, aber Grundideen bis heute in Solvern relevant
\end{itemize}}

\end{frame}

% \begin{frame}\frametitle{DPLL: Überblick}
% 
% DPLL arbeitet mit Formeln in Klauselform (muss vorab gebildet werden)\\[1ex]
% {\footnotesize \phantom{$\leadsto$ }Wir testen Erfüllbarkeit.\\
% $\leadsto$ Hilfsatome können verwendet werden um eine "`kleine"' (polynomielle) Klauselform zu bilden.
% 
% }\bigskip
% 
% Während der Abarbeitung wird schrittweise eine Zuweisung von Wahrheitswerten konstruiert, die bei Erfolg als Modell ausgegeben werden kann.
% % \bigskip
% % 
% % Der Algorithmus verwendet drei wesentliche Funktionen:
% % \begin{itemize}
% % \item \alert{Splitting:} Teile die Suche entsprechend der möglichen Belegungen eines Atoms
% % \item \alert{Unit propagation:} Ermittle einfache Schlussfolgerungen mit Resolution
% % \item \alert{Pure literal elimination:} Entferne trivial erfüllbare Klauseln
% % \end{itemize}
% 
% \end{frame}

\begin{frame}\frametitle{Rückblick: Vereinfachung von $\top$ und $\bot$}

Mit den Äquivalenzen aus Vorlesung 2 erhalten wir folgende Vereinfachungen (anwendbar auf beliebige Teilformeln):
% 
\begin{align*}
%
\top\wedge F  \equiv F\wedge \top & \equiv F		& F\to \top & \equiv \top\\
\top \vee F\equiv F\vee \top & \equiv \top    & \top\to F & \equiv F\\[1ex]
%
\bot\wedge F\equiv F\wedge \bot & \equiv \bot	& F\to \bot & \equiv \neg F\\
\bot\vee F\equiv F\vee \bot & \equiv F		& \bot\to F & \equiv \top\\[1ex]
%
\neg\top & \equiv \bot	& \top\leftrightarrow F\equiv F\leftrightarrow\top & \equiv F\\
\neg\bot & \equiv \top	& \bot\leftrightarrow F\equiv F\leftrightarrow\bot & \equiv \neg F
\end{align*}

Zusammen mit den weiteren Äquivalenzen aus Vorlesung 2 kann jede Formel soweit vereinfacht werden, dass sie entweder $\top$ oder $\bot$ ist,
oder keine Vorkommen von $\top$ oder $\bot$ enthält.

\end{frame}

\begin{frame}\frametitle{Splitting}

\emph{Idee:} Weise einem Atom $p$ einen Wert zu und ermittle (rekursiv) ob die Formel damit erfüllbar ist
\bigskip\pause

Daraus ergibt sich auch schon für beliebigen Formeln ein möglicher Algorithmus:
\codebox{\emph{SplitSat}\\
Eingabe: Formel $F$ (in beliebiger Form)\\
Ausgabe: "`erfüllbar"' oder "`unerfüllbar"'\\[1ex]

\begin{enumerate}[(1)]
\item Wende die Äquivalenzen zur Vereinfachung von $\top$ und $\bot$ erschöpfend auf $F$ an
\item Wenn $F=\bot$: gib "`unerfüllbar"' aus
\item Wenn $F=\top$: gib "`erfüllbar"' aus
\item Andernfalls wähle ein Atom $p$ in $F$:
\begin{itemize}
\item Wenn \emph{SplitSat}($F[p\mapsto\top]$)="`erfüllbar"', dann gib "`erfüllbar"' aus
\item Andernfalls gib \emph{SplitSat}($F[p\mapsto\bot]$) aus
\end{itemize}
\end{enumerate}
}

\emph{Anmerkung:} $F[p\mapsto\top]$ ist die Formel, die aus $F$ durch Ersetzung aller Vorkommen von $p$ durch $\top$ entsteht; analog für $F[p\mapsto\bot]$\\

\end{frame}

\begin{frame}\frametitle{Splitting: Beispiel}

% Beispiel: $(A\leftrightarrow\neg B)\wedge(B\leftrightarrow\neg C)\wedge(C\leftrightarrow(\neg A\wedge\neg B))$
% \bigskip

\begin{forest}
for tree={
%     font=\ttfamily,
    grow'=0,
    child anchor=west,
    parent anchor=south,
    anchor=west,
    calign=first,
    outer ysep=-0.2mm,
    edge path={
      \noexpand\path [draw, \forestoption{edge}]
      (!u.south west) +(7.5pt,0) |-  (.child anchor)\forestoption{edge label};
    },
    before typesetting nodes={
      if n=1
        {insert before={[X,grow'=0,child anchor=west, parent anchor=south, anchor=west, fit=band, calign=first, phantom]}}
        {}
    },
    fit=band,
    before computing xy={l=40pt},
  }
[$(A\leftrightarrow\neg B)\wedge(B\leftrightarrow\neg C)\wedge(C\leftrightarrow(\neg A\wedge\neg B))$
	[$(\top\leftrightarrow\neg B)\wedge(B\leftrightarrow\neg C)\wedge(C\leftrightarrow(\neg\top\wedge\neg B))$, vel={$A\mapsto\top$}, visible on=<2->]
	[$\neg B\wedge(B\leftrightarrow\neg C)\wedge \neg C$,eqnode, visible on=<3->
		[$\neg\top\wedge(\top\leftrightarrow\neg C)\wedge \neg C \equiv \bot$, vel={$B\mapsto\top$}, visible on=<4-> ]
% 		[$\bot$, eqnode]
		[$\neg\bot\wedge(\bot\leftrightarrow\neg C)\wedge \neg C \equiv C\wedge\neg C$, vel={$B\mapsto\bot$}, visible on=<5->
			[ $\top\wedge\neg\top\equiv\bot$, vel={$C\mapsto\top$}, visible on=<6->]
			[ $\bot\wedge\neg\bot\equiv\bot$, vel={$C\mapsto\bot$}, visible on=<7->]
		]
	]
	[$(\bot\leftrightarrow\neg B)\wedge(B\leftrightarrow\neg C)\wedge(C\leftrightarrow(\neg\bot\wedge\neg B))$, vel={$A\mapsto\bot$}, visible on=<8->]
	[$B\wedge (B\leftrightarrow\neg C) \wedge (C\leftrightarrow\neg B)$,eqnode, visible on=<9->
		[$\top\wedge (\top\leftrightarrow\neg C) \wedge (C\leftrightarrow\neg \top)\equiv \neg C$, vel={$B\mapsto\top$}, visible on=<10->
			[$\neg\top\equiv \bot$, vel={$C\mapsto\top$}, visible on=<11->]
			[$\neg\bot\equiv \top$, vel={$C\mapsto\bot$}, visible on=<12->]
		]
	]
]
\end{forest}

\end{frame}

\begin{frame}\frametitle{Splitting: Beobachtungen}

\emph{SplitSat auf allgemeinen Formeln:}
\begin{itemize}
\item rekursive Implementierung der systematischen Betrachtung aller Belegungen\\
	$\leadsto$ vollständig, korrekt, möglicherweise ineffizient
\item Reihenfolge der Belegung der Atome kann Performance beeinflussen
\item Backtracking in Rekursion eingebaut; kann als Suchbaum illustriert/vorgestellt werden
\end{itemize}\bigskip\pause

\emph{Splitting in DPLL:}
\begin{itemize}
\item Bei Formeln in Klauselform genügt eine einzige Form der Vereinfachung: \\\alert{Unit propagation}
\item Weitere sinnvolle Optimierung: \alert{Pure literal elimination}
\end{itemize}

\end{frame}

\begin{frame}\frametitle{Unit Propagation}

\emph{Idee:} Klauseln mit genau einem Literal ("`Unit Clauses"') legen den Wahrheitswert ihres Atoms fest
\begin{itemize}
\item Wir können diesen Wert immer sofort festlegen
\item Alle anderen Klauseln können entsprechend vereinfacht werden ($\leadsto$ Resolution)
\end{itemize}
\bigskip\pause

Für ein Atom $A$ definieren wir $\overline{A}=\neg A$ und $\overline{\neg A}=A$.\smallskip

\codebox{\emph{Unit Propagation}\\
Eingabe: Formel $\mathcal{F}$ in Klauselform\\
Ausgabe: Vereinfachte $\mathcal{F}$ ohne einelementige Klauseln\\[1ex]

Solange es eine Klausel $\{L\}\in\mathcal{F}$ gibt:
\begin{enumerate}[(1)]
% \item Setze den Wert von $L$ auf $\mytrue$ (durch passende Belegung des Atoms in $L$).
\item Lösche alle Klauseln $K\in\mathcal{F}$ mit $L\in K$.
\item Entferne $\overline{L}$ aus allen Klauseln in $K\in\mathcal{F}$ mit $\overline{L}\in K$.
\end{enumerate}
}

\end{frame}

\begin{frame}\frametitle{Pure Literal Elimination}

\emph{Idee:} Atome, die entweder nur nicht-negiert oder nur negiert vorkommen ("`Pure Literals"') sind ohne Probleme erfüllbar
\begin{itemize}
\item Wir können den Wert dieser Atome immer sofort festlegen
\item Alle anderen Klauseln können entsprechend vereinfacht werden ($\leadsto$ Resolution)
\end{itemize}
\bigskip\pause

\codebox{\emph{Pure Literal Elimination}\\
Eingabe: Formel $\mathcal{F}$ in Klauselform\\
Ausgabe: Vereinfachte $\mathcal{F}$ ohne Pure Literals\\[1ex]

Solange in $\mathcal{F}$ ein Pure Literal $L$ vorkommt:
\begin{enumerate}[(1)]
% \item Setze den Wert von $L$ auf $\mytrue$ (durch passende Belegung des Atoms in $L$).
\item Lösche alle Klauseln $K\in\mathcal{F}$ mit $L\in K$.
\end{enumerate}
}
% \medskip

Äquivalent: "`Füge Klausel $\{L\}$ ein und wende Unit Propagation an."'

\end{frame}

\begin{frame}\frametitle{DPLL: Vollständiger Algorithmus}

\codebox{\emph{DPLLSat}\\
Eingabe: Formel $\mathcal{F}$ in Klauselform\\
Ausgabe: "`erfüllbar"' oder "`unerfüllbar"'\\[1ex]

\begin{enumerate}[(1)]
\item Wende Unit Propagation auf $\mathcal{F}$ an
\item Wende Pure Literal Elimination auf $\mathcal{F}$ an
\item Wenn $\mathcal{F}=\emptyset$: gib "`erfüllbar"' aus
\item Wenn $\mathcal{F}$ die leere Klausel enthält: gib "`unerfüllbar"' aus
\item Andernfalls wähle ein Atom $p$ in $\mathcal{F}$:
\begin{itemize}
\item Wenn \emph{DPLLSat}($\mathcal{F}\cup\{p\}$)="`erfüllbar"', dann gib "`erfüllbar"' aus
\item Andernfalls gib \emph{DPLLSat}($\mathcal{F}\cup\{\neg p\}$) aus
\end{itemize}
\end{enumerate}
}

\emph{Anmerkung:} bei Ausgabe "`erfüllbar"' ergeben die Werte der nacheinander in Unit Propagation und Pure Literal Elimination eliminierten Atome ein Modell

\end{frame}

\begin{frame}\frametitle{DPLL: Beispiel}

\begin{forest}
for tree={
%     font=\ttfamily,
    grow'=0,
    child anchor=west,
    parent anchor=south,
    anchor=west,
    calign=first,
    outer ysep=-0.2mm,
    edge path={
      \noexpand\path [draw, \forestoption{edge}]
      (!u.south west) +(7.5pt,0) |-  (.child anchor)\forestoption{edge label};
    },
    before typesetting nodes={
      if n=1
        {insert before={[X,grow'=0,child anchor=west, parent anchor=south, anchor=west, fit=band, calign=first, phantom]}}
        {}
    },
    fit=band,
    before computing xy={l=40pt},
  }
[{$\{\{\neg A, \neg B\}, \{A, B\}, \{\neg B, \neg C\},\{B, C\}, \{\neg C, \neg A\},\{A, B, C\}\}$}
	[{$\{\{A\},\{\neg A, \neg B\}, \{A, B\}, \{\neg B, \neg C\},\{B, C\}, \{\neg C, \neg A\},\{A, B, C\}\}$}, vel={$A\mapsto\top$}, visible on=<2->]
	[{$\{\{\neg B\}, \{\neg B, \neg C\},\{B, C\}, \{\neg C\}\}$}, eqnode, visible on=<3->]
	[{$\{\{C\}, \{\neg C\}\}$}, eqnode, visible on=<4->]
	[{$\{\{\}\}=\{\bot\}$}, eqnode, visible on=<5->]
	%
	[{$\{\{\neg A\},\{\neg A, \neg B\}, \{A, B\}, \{\neg B, \neg C\},\{B, C\}, \{\neg C, \neg A\},\{A, B, C\}\}$}, vel={$A\mapsto\bot$}, visible on=<6->]
	[{$\{\{B\}, \{\neg B, \neg C\},\{B, C\},\{A, B, C\}\}$}, eqnode, visible on=<7->]
	[{$\{\{\neg C\}\}$}, eqnode, visible on=<8->]
	[{$\{\}=\top$}, eqnode, visible on=<9->]
]
\end{forest}

\end{frame}

\begin{frame}\frametitle{DPLL terminiert}

\theobox{\emph{Satz:} Der Algorithmus DPLLSat terminiert für jede Eingabe.}\pause

\emph{Beweis:}
Die Zahl der Atome in $\mathcal{F}$ verringert sich in jedem Durchlauf von Unit Propagation oder Pure Literal Elimination.\\
$\leadsto$ die Schleifen dieser Funktionen terminieren\medskip\pause

Rekursive Aufrufe sind nur unter Hinzunahme von Unit Clauses möglich\\
$\leadsto$ mit jeder Rekursion sinkt die Zahl der Atome mindestens um $1$\\\pause
$\leadsto$ die maximale Rekursionstiefe ist linear
\qed
\bigskip

Aus der bestimmten maximalen Rekursionstiefe folgt:\\
\alert{DPLLSat führt zu maximal $2^k$ rekursiven Aufrufen ($k$ = Zahl der Atome in $\mathcal{F}$).}

\end{frame}


\begin{frame}\frametitle{DPLL ist korrekt}

\theobox{\emph{Satz:} Das Ergebnis von DPLLSat($\mathcal{F}$) ist korrekt.}\pause

\emph{Beweis:}
Die Vereinfachungen produzieren äquivalente Formeln (unter den gewählten Belegungen), weil das Löschen von Klauseln oder Literalen
Sonderformen der bereits bekannten Äquivalenumformungen sind.
\medskip\pause

Korrektheit folgt durch Induktion über Zahl $k$ der Atome in $\mathcal{F}$ den Vereinfachungen\pause:
\begin{itemize}
	\item Für $k=0$ Atome ist die Klauselform $\{\}$ oder $\{\{\}\}$. In beiden Fällen ist das Ergebnis korrekt.\pause
	\item Bei $k>0$ Atomen führt Splitting zu rekursiven Aufrufen für $\mathcal{F}\cup\{p\}$ und $\mathcal{F}\cup\{\neg p\}$.\pause
	In beiden Fällen wird sich die Zahl der Atome durch Unit Propagation um mindestens eins verringern.\pause
	Per Induktionsannahme ist das Ergebnis der rekursiven Aufrufe also korrekt. \pause
	Korrektheit für $k$ folgt, weil $\mathcal{F}$ genau dann erfüllbar ist, wenn $\mathcal{F}\cup\{p\}$ oder $\mathcal{F}\cup\{\neg p\}$
	erfüllbar sind.\qed
\end{itemize}

\end{frame}


\begin{frame}\frametitle{DPLL in der Praxis}

\emph{DPLL ist bereits ein praktikables Verfahren}
\begin{itemize}
\item oft wenige Splittings nötig, speziell bei (leicht) erfüllbaren Formeln
\item Eingaben mit vielen Horn-Formeln fast komplett mit Unit Propagation lösbar
\item insgesamt viel effizienter als Resolution
\end{itemize}
\medskip\pause

\emph{Moderne Solver verwenden weitere Verbesserungen:}
\begin{itemize}
\item heuristische \alert{Strategien zur Wahl der Atome} für das Splitting
\item \alert{Backjumping:} Rückgängigmachung vieler Wertzuweisungen in einem Schritt, überspringt "`irrelevante"' Entscheidungen im Backtracking
\item \alert{Conflict-driven clause learning (CDCL):} Erkenntnisse aus Fehlschlägen in der Suche werden in zusätzlichen Klauseln "`abgespeichert"' um die weitere Suche abzukürzen
\end{itemize}

\end{frame}

\sectionSlide{SMT-Solving}

\begin{frame}\frametitle{Satisfiability Modulo Theories}

\emph{Aussagenlogisches Problem:}
\begin{center}
% Ist $(\neg p\vee q)\wedge  $
Gibt es Werte für $p,q,r,s,t$ aus der Menge \{\mytrue,\myfalse\}, so dass gilt\\
$(p \vee r)\wedge (q \vee s)\wedge (p\vee q)\wedge \neg t$?
\end{center}\pause

\emph{Mathematisches Problem mit aussagenlogischen Operatoren:}
\begin{center}
% Ist $(x\leq y) \wedge (y\leq z) \wedge ((x<y) \vee (y<z)) \wedge \neg (x<z)$
Gibt es Werte für $x,y,z$ aus der Menge der natürlichen Zahlen, so dass gilt\\
$(x<y \vee x=y)\wedge (y<z \vee y=z)\wedge (x<y\vee y<z)\wedge \neg(x<z)$?
\end{center}
\pause

\alert{Satisfiability Modulo Theories (SMT):}
\begin{itemize}
\item Verallgemeinerung von Aussagenlogik mit Atomen, die sich auf andere (nicht-boolesche) Domänen beziehen
\item typische Datentypen: ganze Zahlen, reelle Zahlen, Listen, Strings, Enums usw.
\item Operationen je nach Typ: Arithmetik, Stringoperationen, Gleicheiten und Ungleichheiten usw.
\end{itemize}
\emph{Ziel:} Ermittle Erfüllbarkeit und finde gegebenenfalls ein Modell

\end{frame}

\begin{frame}[fragile]\frametitle{Beispiel: Z3 als SMT-Solver}

Die Erfüllbarkeit von $(x<y \vee x=y)\wedge (y<z \vee y=z)\wedge (x<y\vee y<z)\wedge \neg(x<z)$ kann in Z3 wie folgt getestet werden:
\medskip

\begin{tikzpicture}%
\node [draw=black!50, fill=white, very thick, rectangle, rounded corners, inner sep=10pt, inner ysep=10pt, align=left, text width=0.9\linewidth] (box){%
\vspace{-2ex}
\begin{minted}{python}
from z3 import *

x,y,z = Ints('x y z')

s = Solver()
s.add(And( Or(x<y, x==y), Or(y<z, y==z), Or(x<y, y<z), Not(x<z)))

print(s.check())
\end{minted}
};%
\end{tikzpicture}

\end{frame}

\begin{frame}\frametitle{SMT und DPLL}

Ein wichtiger Ansatz zum Lösen von SMT-Problemen\\
\alert{(am Beispiel $x\leq y\wedge(x>y\vee x=y)$)}
\begin{itemize}
\item Abstrahiere die Eingabe in eine aussagenlogische Formel\\
\alert{Beispiel: $x\leq y\wedge(x>y\vee x=y)$ wird zu $p\wedge(q\vee r)$}
\item Suche mit DPLL nach erfüllenden Belegungen\\
\alert{Beispiel: $p\mapsto\mytrue, q\mapsto\mytrue, r\mapsto\myfalse$}
\item Übersetze (partielle) Belegungungen zurück in die Ausgangssprache und prüfe dort ihre Erfüllbarkeit
$\leadsto$ benötigt spezifische Verfahren\\
\alert{Beispiel: $x\leq y, x>y, x\neq y$ unerfüllbar}
\item Falls erfüllbar: gib Modell aus
\item Falls unerfüllbar: lerne zusätzliche logische Bedingungen\\
\alert{Beispiel: Formel wird um Klausel $(\neg p\vee\neg q)$ erweitert}
\end{itemize}

\end{frame}


\begin{frame}\frametitle{Das T in SMT steht für Theory}

\alert{Theory} bezeichnet die logisch-mathematische Theorie für einen Datentyp und eine Menge unterstützter Operationen
\medskip

Beispiele:
\begin{itemize}
\item Relle Zahlen mit $*$, $+$, $-$, $/$
\item Ganze Zahlen mit lineraren Gleichungssysteme
\item Strings mit Verkettungsoperationen und lexikographischer Ordnung
\item IEEE Gleitkommazahlen mit Gleitkomma-Arithmetik
\item \ldots
\end{itemize}
\medskip\pause

\emph{Das Schließen in einer Theorie kann sehr kompliziert sein}
\begin{itemize}
\item Jede Theorie verlangt eigene Methoden
\item Viele Theorien sind unentscheidbar $\leadsto$ unvollständige Erfüllbarkeitstests möglich, aber keine Erfolgsgarantie 
\item Die Vermischung von Theorien in einer Formel ist nicht trivial 
\end{itemize}

\end{frame}

\begin{frame}\frametitle{SMT-Solving: Ausblick}

Die "`erweiterten aussagenlogischen Formeln"' in SMT gehören schon in den Bereich der Prädikatenlogik
$\leadsto$ später mehr dazu
\bigskip

"`DPLL mit Theorien"' ist bei weitem nicht die einzige SMT-Methode, aber eine der wichtigsten
\bigskip

SMT-Solver wie Z3 benötigen weitere Techniken, unter anderem
\begin{itemize}
\item Normalisierung und Modifikation von logischen Ausdrücken
\item Vereinfachungen und Äquivalenzumformungen für Ausdrücke aller Theorien
\end{itemize}
\bigskip

Z3 (und andere SMT-Solver) in der Praxis:
\begin{itemize}
\item vielfältig anwendbar, auch ohne detaillierte Kenntnisse
\item oft viele verschiedene Problemkodierungen möglich
\item Performance in der Regel nicht theoretisch vorhersagbar (aber Hintergrundkenntnisse und Erfahrung helfen)
\end{itemize}


\end{frame}

\sectionSlide{Logisches Schließen als Wortproblem}

\begin{frame}\frametitle{Schließen als Entscheidungsproblem}

Wir wissen bereits, dass Schließen ein Entscheidungsproblem ist,\ghost{ z.B.:}

\defbox{\redalert{Erfüllbarkeit}:
\begin{itemize}
\item \emph{Eingabe:} Formel $F$
\item \emph{Ausgabe:} "`ja"' wenn $F$ erfüllbar ist; sonst "`nein."'
\end{itemize}
}

Wie wir wissen, können solche Probleme als Definition einer Sprache angesehen werden:
\[\Slang{SAT} = \{ \textsf{enc}(F)\mid F\text{ ist erfüllbar}\}\]
wobei $\textsf{enc}(F)$ eine (vernünftige) Kodierung der Formel $F$ als Wort über einem endlichen Alphabet (z.B.\ $\{\Sterm{0},\Sterm{1}\}$) ist. Vor allem müssen wir beliebig viele Atome mit nur endlich vielen Zeichen kodieren.
\bigskip\pause

\anybox{purple}{%
\narrowcentering{Was für eine Sprache ist $\Slang{SAT}$?}\\
\narrowcentering{Durch welches Automatenmodell wird sie erkannt?}}

\end{frame}

\begin{frame}[t]\frametitle{Schließen als Sprache (1)}

\anybox{purple}{%
\narrowcentering{Was für eine Sprache ist $\Slang{SAT}$?}\\
\narrowcentering{Durch welches Automatenmodell wird sie erkannt?}}\medskip

Wir wissen, dass $\Slang{SAT}$ entscheidbar ist und daher von Typ 0 ist.\pause\\
Der folgende naive Algorithmus entscheidet $\Slang{SAT}$ auf einer TM:
% Auch für Typ~1 gibt es ein positives Resultat:

\codebox{%Der folgende Algorithmus entscheidet $\Slang{SAT}$ auf einem LBA:
Naiver Erfüllbarkeitstest (Skizze):
\begin{itemize}
\item Wir prüfen systematisch jede Wertzuweisung auf den Atomen der gegebenen Formel
\item Dazu iteriert die TM über alle Wertzuweisungen und berechnet für jede rekursiv den Wahrheitswert der Formel
\item Falls eine erfüllende Zuweisung gefunden wird, dann hält die TM und akzeptiert
\item Falls alle Wertzuweisungen ohne Erfolg durchlaufen worden sind, dann hält die TM und verwirft
\end{itemize}
}

\end{frame}

\newcommand{\vtriple}[3]{\small\begin{pmatrix}#1\\[-1.3ex]#2\\[-1.3ex]\Sterm{#3}\end{pmatrix}}
\newcommand{\vtripletext}[3]{\footnotesize\begin{array}{r}\text{#1}\\[-1.0ex]\text{#2}\\[-1.0ex]\text{#3}\end{array}}

\begin{frame}[t]\frametitle{Schließen als Sprache (2)}

\anybox{purple}{%
\narrowcentering{Was für eine Sprache ist $\Slang{SAT}$?}\\
\narrowcentering{Durch welches Automatenmodell wird sie erkannt?}}\smallskip

Der naive Erfüllbarkeitstest kann in linearem Speicher ablaufen, d.h. auf einem LBA \alert{$\leadsto$ $\Slang{SAT}$ ist von Typ 1}\pause
% Auch für Typ~1 gibt es ein positives Resultat:

\codebox{%
Naive Erfüllbarkeit auf LBA (Skizze):\footnotesize
\begin{itemize}%
\item Wir verwenden ein Arbeitsalphabet mit mehreren "`Spuren"': (1)~Eingabe, (2)~aktuell ermittelte Wahrheitswerte aller Teilformeln, (3)~aktuell getestete Wertzuweisung\vspace{-0.5ex}
\item Skizze der Kodierung (ohne Binärkodierung der Formel):\\[1ex]
\narrowcentering{$
\vtripletext{Spur 3}{Spur 2}{Spur 1}
% 
\vtriple{}{}{(}
\vtriple{}{}{(}
\vtriple{\mytrue}{\mytrue}{p}
\vtriple{}{\myfalse}{\to}
\vtriple{\myfalse}{\myfalse}{q}
\vtriple{}{}{)}
\vtriple{}{\mytrue}{\vee}
\vtriple{}{}{(}
\vtriple{}{\myfalse}{q}
\vtriple{}{\mytrue}{\to}
\vtriple{\myfalse}{\myfalse}{r}
\vtriple{}{}{)}
\vtriple{}{}{)}
$}
\item Man kann nun systematisch Spur 3 iterieren, Spur 2 rekursiv berechnen und den Wert der Gesamtformel prüfen
(machbar, wir verzichten auf die Details)
% \item Dabei ordnen wir jeder Teilformel eine eindeutige Position zu: bei Atomen ist dies das
% erste Zeichen der Kodierung des Atoms, bei zusammengesetzten Formeln das erste Zeichen der Kodierung des Junktors
% \item Der aktuelle Wahrheitswert der Teilformel wir in Spur~2 an ihrer Position gespeichert.
% \item Die aktuelle Wertzuweisung ergibt sich aus den Wahrheitswerten der Atome
% \item Wenn man das erste Vorkommen jedes Atoms in der Formel markiert, dann ist es nicht schwer, die jeweils nöchste Wertzuweisung zu berechnen (durch binäre Addition von $1$).
\end{itemize}\vspace{-3.5ex}~}

\end{frame}

\begin{frame}\frametitle{Welches Berechnungsmodell passt zu $\Slang{SAT}$?}

Erkenntnisse bisher:
\begin{itemize}
\item LBAs sind stark genug für $\Slang{SAT}$ (d.h. $\Slang{SAT}$ ist von Typ 1)
\item $\Slang{SAT}$ ist keine reguläre Sprache, da FAs korrekte Klammerung nicht erkennen können (d.h. $\Slang{SAT}$ ist nicht von Typ 3)
% \item ??? $\Slang{SAT}$ ist keine kontextfreie Sprache (d.h. $\Slang{SAT}$ ist nicht von Typ 2)
\end{itemize}\bigskip\pause

Man kann außerdem zeigen:
\theobox{\narrowcentering{$\Slang{SAT}$ ist nicht kontextfrei.}}
(zumindest wenn eine offensichtliche Kodierung verwendet wird; siehe
\url{http://cstheory.stackexchange.com/questions/37322/is-sat-a-context-free-language})
\bigskip\pause

\anybox{purple}{%
Können wir $\Slang{SAT}$ noch genauer charakterisieren als mit LBAs?}\\\pause
$\leadsto$ Ja, mithilfe anderer Arten von eingeschränkten TMs.

\end{frame}


\begin{frame}\frametitle{Zusammenfassung und Ausblick}

\redalert{DPLL} ist ein effizientes Verfahren zur Ermittlung der Erfüllbarkeit aussagenlogischer Formeln,
basierend auf systematischer Modellsuche
\bigskip

Wichtige darin vorkommende Methoden sind \redalert{Backtracking} (und Splitting), \redalert{Unit Propagation} und
die Eliminierung von \redalert{Pure Literals}\bigskip

DPLL kann auf nicht-boolesche Domänen erweitert werden, was nützlich für SMT-Solving ist, aber Resolution ist besser
auf ausdrucksstarke Logiken übertragbar (später mehr dazu)
\bigskip

\anybox{yellow}{
Offene Fragen:
\begin{itemize}
% \item Welche Methoden des logischen Schließens sind praktikabel?
% \item Ist Resolution korrekt?
% \item Ist Resolution effizient?
% \item Ist logisches Schließen immer so schwer?
% \item Gibt es auch sub-exponentielle Algorithmen für aussagenlogisches Schließen?
\item Was hat Aussagenlogik mit Sprachen, Berechnung und TMs zu tun?
\item Unsere Ansätze zum logischen Schließen sind alle schlimmstenfalls exponentiell -- muss das sein?
\end{itemize}
}
\end{frame}

\end{document}
